\documentclass[titlepage]{article}
\usepackage{hyperref}
\usepackage{amsthm}
\usepackage{mathrsfs}
\usepackage[margin=1in]{geometry}
\usepackage{amssymb}
\usepackage{amsmath}
\usepackage{tikz-cd}

\newcommand{\Exercise}[1]{\textbf{Exercise #1}}
\newcommand{\Solution}{\textbf{Solution }}
\newcommand{\cat}[1]{\mathsf{#1}}
\newcommand{\Cat}{\cat{C}}
\newcommand{\id}{\text{id}}
\DeclareMathOperator{\End}{End}
\DeclareMathOperator{\Hom}{Hom}
\DeclareMathOperator{\Obj}{Obj}



\title{Algebra: Chapter 0 Solutions}
\author{Dylan Huis in't Veld}
\begin{document}
\maketitle
\tableofcontents

\newpage

\section*{Preface}

This document does not seek to represent the perfect solutions manual to the text at hand, and is more of a way of recording my progress through the text. Should you happen to find this, and wish to contribute or point out any logical flaws in the arguments presented, I would be more than happy to adjust my solutions.

My main goal of writing this is to improve my \LaTeX fluency, having previously gotten used to Typst. However, having grown interested in algebra after taking courses in group theory and topology (which included elements of algebraic topology) at university, I wanted to take a more in depth route through the subject while also exploring further structures beyond group theory. I apologise if my solutions to the group theory questions are not as thorough, as this will come subconsciously as a result of likely having done the exercises, or those similar to, before.

My progress through this will likely decrease once the semester starts back up, as I have started this endeavour during my winter break. I apologise if you are waiting for solutions for a while. I will do my best to work on this when I can.
\newpage

\section{Chapter 1: Preliminaries: Set Theory and Categories}
\subsection{Naive Set Theory}

\textbf{Exercise 1.1.} Locate a discussion of Russell's paradox, and understand it.

\medskip

\textbf{Solution} There are many discussions of the paradox. The one I found most intuitive or easy to understand was the labelling of sets as such: We say a set is 'abnormal' if the set contains itself ($S \in S$), or we say that a set is 'normal' if it does not contain itself ($S \notin S$). We now define the set $X$ as the set of all normal sets. We ask the question 'is $X$ a normal or abnormal set?'

\medskip

 Suppose it is normal, then it would be contained within $X$ by definition, making $X$ abnormal. If we instead suppose $X$ is abnormal, then $X$ would not be contained in itself, making $X$ normal.


\medskip
 We arrive at a paradox; that is, $X$ is neither normal or abnormal.

 \qed

\bigskip

\textbf{Exercise 1.2.} Prove that if $\sim$ is an  equivalence relation on a set $S$, then the corresponding family $\mathscr{P}_{\sim}$ defined in \S 1.5 in indeed a partition of $S$: that is, its elements are nonempty, disjoint and their union is $S$. 

\medskip

\textbf{Solution} We show the desired result by considering the definition of the equivalence class of $a$ given in the text

\[ [a]_{\sim} := \{b \in S | b \sim a\} \]

Firstly, its clear that every element $s \in S$ is contained in its own equivalence class. This is due to the reflexivity requirement of the relation ($s \sim s$). This also shows that every equivalence class is non-empty. We now show that for any two equivalence classes $[a]_{\sim}$ and $[b]_{\sim}$ are either disjoint or equal.

Suppose the two sets are not disjoint. Then there exists an element $z \in [a]_{\sim} \cap [b]_{\sim}$. Let $n \in [a]_{\sim}$ be arbitrary. Since we have $n \sim z$ and $z \sim b$, we must have that $n \in [b]_{\sim}$ by transitivity and so $[a]_{\sim} \subseteq [b]_{\sim}$. Similarly, for an arbitrary element $m \in [b]_{\sim}$ we have $m \sim z$ and $z \sim a$ and so $m \in [a]_{\sim}$ implying that $[b]_{\sim} \subseteq [a]_{\sim}$. Combining these two results means that $[a]_{\sim} = [b]_{\sim}$ and the two equivalence classes are equal.

Totalling the results thus far, we see that all elements of $S$ are in a equivalence class, and that these equivalence classes are either disjoint or equal, and so $\mathscr{P}_{\sim}$ indeed forms a partition of $S$.

\qed

\bigskip

\textbf{Exercise 1.3.} Given a partition $\mathscr{P}$ on a set $S$, show how to define a relation $\sim$ on $S$ such that $\mathscr{P}$ is the corresponding partition.

\medskip

\textbf{Solution} Let $P_i$ denote the disjoint non empty subsets in $\mathscr{P}$. We define an equivalence relation on elements of $S$ by the following: $a \sim b$ if and only if $a$ and $b$ are contained within the same subset $P_i$. We need only show this defines a valid equivalence relation.

We clearly have that $a \sim a$ vacuously by the definition of the equivalence relation, so the reflexivity requirement is satisfied. Supposing that $a \sim b$, this means that $a,b \in P_i$ and so equivalently we have that $b \sim a$, satisfying symmetry. Suppose on top of this we have that $b \sim c$. Then this must mean that $c$ is also the in same subset $P_i$, thereby meaning that $a \sim c$, satisfying transitivity, making this a well defined equivalence relation. By construction this relation generates the exact same partition of $S$ that $\mathscr{P}$ does.

\qed 

\bigskip

\textbf{Exercise 1.4.} How many different equivalence relations may be defined on the set $\{1,2,3\}$.

\medskip

\textbf{Solution} Combining the results of the two previous exercises we know that an equivalence relation defines a partition, $\mathscr{P}_{\sim}$ of a set $S$ by equivalence classes, and that given any partition $\mathscr{P}$ of $S$, we can define a corresponding equivalence relation that defines this partition. What this means is that this question is as simple as counting all partitions of the set. We have the followings partitions of $\{1,2,3\}$:

\[ \{ \{1,2,3\} \}, \{ \{1,2 \}, \{3 \} \}, \{ \{1 \}, \{2,3 \} \}, \{ \{1,3 \}, \{2 \} \}, \{ \{1 \}, \{2 \}, \{3 \} \} \]

Giving a total of 5 equivalence relations.

\qed

\bigskip

\textbf{Exercise 1.5.} Give an example of a relation that is reflexive and symmetric but not transitive. What happens if you attempt to use this relation to define a partition on the set?(Hint: Thinking about the second question will help you answer the first one.)

\medskip

\textbf{Solution} Rather than coming up with a strict rule for the relation, we will simply construct a family of subsets that is not a partition and show how this defines the desired relation. Take $\{1,2,3 \}$ and consider the following family of subsets:

\[ \{ \{1 \}, \{2 \}, \{3 \}, \{1,2 \}, \{2,3 \} \} \]

If we define this as equivalence classes, we get that the relation is reflexive due to the first 3 subset in the family. We also have that $1 \sim 2 \implies 2 \sim 1$ since there is a subset containing both elements, and the same is true for $2 \sim 3$ satisfying the symmetry requirement. However we do not have $1 \sim 3$ despite having $1 \sim 2$ and $2 \sim 3$, and so transitivity is not fulfilled. The main observation to take away here is that if whatever relation we are working with fails to satisfy all 3 requirements, it will result in a family of subsets that are not a partition of the original set, and vice versa.

\qed

\bigskip

\textbf{Exercise 1.6.} Define a relation $\sim$ on the set $\mathbb{R}$ of real numbers by setting $a \sim b \iff b - a \in \mathbb{Z}$. Prove that this is an equivalence relation, and find a 'compelling' description for $\mathbb{R} / \sim$. Do the same for the relation $\approx$ on the plane $\mathbb{R} \times \mathbb{R}$ defined by declaring $(a_1,a_2) \approx (b_1,b_2) \iff b_1 - a_1 \in Z$ and $b_2 - a_2 \in Z$.

\medskip

\textbf{Solution} For any $a \in \mathbb{R}$ we have that $a-a=0 \in \mathbb{Z}$ so $a \sim a$, satisfying reflexivity. If $a \sim b$ then $a-b \in \mathbb{Z}$ which implies that $-(a-b) = b-a \in Z$ and so $b \sim a$, satisfying symmetry. If $a \sim b$ and $b \sim c$ then we have that $a-b$ and $b-c$ are both integers. We must then have that $(a-b) + (b-c)$ is also an integer, and this is exactly equal to $a-c$, and so we then have that $a \sim c$, satisfying transitivity and thereby proving this is a valid equivalence relation. The second equivalence relation given in the exercise is exactly the same, just working element wise on each ordered pair, with the workings being analogous to those done already.
\smallskip

In terms of a description, we can think of the first equivalence relation as the real number line being quotiented down to just the interval $[0,1]$. Similarly, for the second relation this is a quotient of the cartesian plane where it is quotiented down to the unit square $[0,1]^2$. If one is familiar with topology, they may know that this set endowed with the quotient topology coming from the standard topology on $\mathbb{R}^2$ gives the 2-torus $\mathbb{T}^2$.

\bigskip

\subsection{Functions Between Sets}

\textbf{Exercise 2.1.} How many different bijections are there between a set $S$ with $n$ elements and itself?

\medskip

\textbf{Solution} This can be solved by thinking about the different bijections element-wise in the set. We have $n$ choices for where to send the first element in the set, and once this is chosen we will then have $(n-1)$ choices of where to send our second element. We can continue this process to find the the total number of bijections is equal to 

\[ n \times (n-1) \times (n-2) \times \dots 2 \times 1 = n! \]

We know by construction these are indeed injective and surjective, and we will cover every single possible combination of element-wise pairings this way, confirming this is indeed the total number of bijections.

\qed

\bigskip

\textbf{Exercise 2.2.} Prove statement (2) in Proposition 2.1. You may assume that given a family of disjoint nonempty subsets of a set, there is a way to choose one element in each member of the family.

\medskip

\textbf{Solution} The claim we must prove is the following: \textit{$f$ has a right-inverse if and only if it is surjective}.

$(\implies)$ Suppose $f: A \to B$ has a right-inverse $g:B \to A$ so that

\[ f \circ g = \textrm{id} _B \]

Let $b \in B$ be any arbitrary elements. We then have that $(f \circ g)(b) = \textrm{id}_B (b) = b$. This then means that $f(g(b))=b$ and so we have found a suitable element $g(b) \in A$ that $f$ maps to $b$, implying that $f$ is surjective.

$(\impliedby)$ Conversely suppose that $f$ is surjective. For every $b \in B$ we consider the fiber of $f$ over $b$, $f^{-1}(b)$. Since $f$ is not necessarily injective, each fiber can have any number of elements. However, by the definition of a function, we know that each fiber is disjoint, and by the axiom of choice we can choose an $a \in f^{-1}(b)$ from each fiber. We then define a function $g:B \to A$ by setting $g(b)=a$. We then have that
\[ (f \circ g)(b) = f(g(b)) = f(a) = b \]

Implying that $f \circ g = \text{id}_B$ and so $f$ has a right inverse.

\qed

\bigskip

\textbf{Exercise 2.3.} Prove that the inverse of a bijection is a bijection and that the composition of two bijections is a bijection.

\medskip

\textbf{Solution} Let $f:A \to B$ be a bijection, and denote its inverse $f^{-1}:B \to A$

\textit{Injective}: Suppose that $f^{-1}(b') = f^{-1}(b'')$. By injectivity of $f$ we must then have that $f(f^{-1}(b')) = f(f^{-1}(b''))$. By the definition of the inverse, which states that the composition of $f$ and its inverse it equal to the identity function, this then implies that $b' = b''$. Therefore, $f^{-1}$ is injective.

\smallskip

\textit{Surjective} Let $a \in A$. we then have that $f(a) \in B$ and $f^{-1}(f(a)) = a$ so $f^{-1}$ is surjective. We have thus proved that the inverse is also a bijection.

\medskip

Let $g:C \to A$ be another bijection. We will show that $f \circ g$ is a bijection.

\textit{Injective}: Let $c',c'' \in C$ and suppose that $(f \circ g)(c') = (f \circ g)(c'')$. We then have that
\[ f(g(c')) = f(g(c'')) \]

By the injectivity of $f$ we must have that $g(c')=g(c'')$, and then by injectivity of $g$ we must then have that $c' = c''$ showing that $f \circ g$ is injective.

\smallskip

\textit{Surjective}: Let $b \in B$ be an arbitrary element. The surjectivity of $f$ means $\exists a \in A$ so that $f(a) = b$. Then by surjectivity of $g$ $\exists c \in C$ so that $g(c)=a$. We then observe the following

\[ (f \circ g)(c) = f(g(c)) = f(a) = b \]

Showing that $f \circ g$ is also surjective and thus is a bijection.

\qed

\bigskip

\textbf{Exercise 2.4.} Prove that 'isomorphism' is an equivalence relation (on any set of sets).

\smallskip

\textbf{Solution} We show that each requirement of an equivalence relation holds, which comes almost immediately from the previous exercises.

\textit{Reflexive}: For any set $A$ we must have $A \cong A$ as the identity function on $A$, $\text{id}_A:A \to A$ is a bijection from $A$ to $A$ and so $A$ is isomorphic to itself.

\textit{Symmetric}: Suppose that $A \cong B$. This then means there exists a bijection, say $f:A \to B$, between the two sets. Exercise (2.3) proved that the inverse of $f$, $f^{-1}:B \to A$ is also a bijection, and so this then means that $B \cong A$.

\textit{Transitive}: Suppose we have $A \cong B$ and $B \cong C$. This then means we have two bijections $f:A \to B$ and $g:B \to C$. Exercise (2.3) also demonstrated that the composition of bijections is a bijection, and so $g \circ f:A \to C$ is a bijection and we have $A \cong C$.

\qed

\bigskip

\textbf{Exercise 2.5.} Formulate the notion of \textit{epimorphism}, in the style of the notion of \textit{monomorphism} seen in \S 2.6, and prove a results analogous to Proposition 2.3, for epimorphisms and surjection.

\medskip

\textbf{Solution} An epimorphism can be defined as follows: A function $f:A \to B$ is an \textit{epimorphism} if the following holds:

\begin{center}
\textit{for all sets $Z$ and all functions $\alpha ',\alpha '': B \to Z$}

\[ \alpha ' \circ f = \alpha '' \circ f \implies \alpha ' = \alpha '' \]
\end{center}

The analogous proposition we wish to prove is that \textit{A function is surjective if and only if it is an epimorphism}

($\implies$) Suppose $f:A \to B$ is surjective. Then $f$ has a right inverse. For any functions $\alpha ',\alpha '': B \to Z$ we have

\[ \alpha ' \circ f = \alpha '' \circ f \implies \alpha ' \circ f \circ g = \alpha '' \circ f \circ g \implies \alpha ' \circ \text{id}_A = \alpha '' \circ \text{id}_A \implies \alpha ' = \alpha '' \]

so $f$ is an epimorphism.

($\impliedby$) Converse suppose that $f$ is an epimorphism. Then define $\alpha ', \alpha '': B \to \{0,1\}$ by the following:

\[ \alpha '(b) = \begin{cases}
1 & b \in \text{im} f \\ 0 & b \notin \text{im} f
\end{cases}, \alpha '' (b) = 1 \]


Then clearly we have that $\alpha ' \circ f = \alpha '' \circ f$, and by the definition of $f$ being an epimorphism we must have $\alpha ' = \alpha ''$. This can only be true if $B = \text{im} f$ which means that $f$ must be surjective.

\qed

\bigskip

\textbf{Exercise 2.6.} With notation as in Example 2.4, explain how any function $f:A \to B$ determine a section of $\pi_A$.

\medskip

\textbf{Solution} Recall that a section is just a right inverse, so we want to find a function $g:A \to A \times B$ such that

\[ \pi_A \circ g = \text{id}_{A} \]

Since $\pi_A((a,b)) = a$, we want to construct $g$ such that $g(a) = (a,b)$ where b is any arbitrary element of $B$ that is not important to us. However, since we want our section to be defined by $f$ we can set $g(a)=(a,f(a))$. This means g is determined by $f$, and by construction we have that $g$ is a section of $\pi_A$.

\qed

\bigskip

\textbf{Exercise 2.7.} Let $f:A \to B$ be any function. Prove that the graph $\Gamma_f$ of $f$ is isomorphic to $A$.

\medskip

\textbf{Solution} Whenever encountering a question like this, I think its more fruitful to design a function going from the smaller set, $A$, to the bigger set $\Gamma_f$. This is largely due to the fact that when going the other way, you are losing some information, and in cases like this it can be hard to make the function injective. We will construct our function as follows:

\[ g:A \to \Gamma_f, a \mapsto (a,f(a)) \]

We show that this is a bijection.

\textit{Injective}: Suppose that $g(a)=g(a')$. This must then mean that $(a,f(a))=(a',f(a'))$, implying that $a = a'$.

\textit{Surjective}: for any $(a,f(a)) \in \Gamma_f$ we have that $g(a) = (a,f(a))$ and so $g$ is clearly surjective.

We have constructed a bijection between the two sets, and so we have $A \cong \Gamma_f$.

\qed

\bigskip

\textbf{Exercise 2.8.} Describe as explicitly as you can all terms in the canonical decomposition (cf. \S 2.8)
 of the function $\mathbb{R} \to \mathbb{C}$ defined by $r \mapsto e^{2 \pi i r}$. (This exercise matches one assigned previously. Which one?)

\medskip

\textbf{Solution} The exercise this matches is Exercise (1.6). The reason this is the case is that the equivalence relation formed from this mapping creates equivalence classes such that the distance between each element is an integer value. In other words, we we define this map as $g$, then $a' \sim a'' \iff g(a')=g(a'') \iff a'-a'' \in \mathbb{Z}$ Which is exactly the recurrence relation we had in the mentioned exercise. Letting $\sim$ represent the relation mentioned, we can now list every element of our canonical decomposition

\[ s:\mathbb{R} \to \mathbb{R} / \sim, a \mapsto [a]_{\sim} \]

\[ \tilde{g}: \mathbb{R} / \sim \to \text{im} f, [a]_{\sim} \mapsto e^{2 \pi i a} \] 

\[ i: \text{im} f \to \mathbb{C}, c \mapsto c \]

And the decomposition can be realised in the following diagram:

\[
\begin{tikzcd}
\mathbb{R} \arrow[r,twoheadrightarrow, "s" below,] \arrow[rrr,"f",bend left=30] & \mathbb{R} /\sim \arrow[r,"\sim" above,"\tilde{g}" below] &  \text{im} f \arrow[r,hook,"i" below] & \mathbb{C}
\end{tikzcd}
\]

\qed

\bigskip

\textbf{Exercise 2.9} Show that if $A' \cong A''$ and $B' \cong B''$, and further $A' \cap B' = \emptyset$ and $A'' \cap B'' = \emptyset$, then $A' \cup B' \cong A'' \cup B''$. Conclude that the operation $A \amalg B$ (as described in \S 1.4) is well-defined up to \textit{isomorphism} (cf. \S 2.9).

\medskip

\textbf{Solution} We show that this is true by constructing the explicit bijection between the two sets. Since $A' \cong A''$, there exists a bijection $\alpha: A' \to A''$. Similarly, since $B' \cong B''$ there exists a bijection $\beta: B' \to B''$. We construct a new function, $\gamma:A' \cup B' \to A'' \cup B''$ as follows:


\[ \gamma(x) = \begin{cases}
\alpha(x) & x \in A' \\ \beta(x) & x \in B'
\end{cases} \]

Since $A'$ and $B'$ are disjoint, this function is well defined, and its also clear that if we consider $\gamma\vert_{A'}$ and $\gamma\vert_{B'}$, that gamma inherits the bijective nature from each set, and since the sets are disjoint we can be sure that $\gamma$ is bijective across its entire domain. This thereby means that $A' \cup B' \cong A'' \cup B''$.


When we constructed $A \amalg B$, we found two sets $A'$ and $B'$ which were isomorphic to $A$ and $B$ respectively and took the normal union of these two sets. The above computation verifies that no matter which isomorphic sets we choose, the union will be the same up to isomorphism. This is exactly what it means for this operation to be well defined up to isomorphism.

\qed

\bigskip

\textbf{Exercise 2.10.} Show that if $A$ and $B$ are finite sets, then $|B^A|=|B|^{|A|}$

\medskip

\textbf{Solution} This exercise is very similar to (2.1) and we proceed with the same counting argument. Recall that $B^A$ is the set of all functions from a set $A$ to a set $B$. Since we are not concerned about injectivity, there is no restriction on where we map each element of $A$. Therefore, we have $|B|$ options of where to map each individual element of $A$, and since we have $|A|$ elements we need to map, there are exactly $|B|^{|A|}$ unique functions from $A$ to $B$.

\qed

\bigskip

\textbf{Exercise 2.11.} In view of Exercise 2.10, it is not unreasonable to use $2^A$ to denote the set of functions from an arbitrary set $A$ to a set with 2 elements (say {0,1}). Prove that there is a bijection between $2^A$ and the \textit{power set} of A.


\textbf{Solution} Encountering your first question of this nature can be somewhat daunting. I think it helps to take a step back and simplify what we are wanting to show. We want to match any subset of $A$ to a function $f:A \to \{0,1\}$ such that we guarantee there is no overlap, and every function is covered. Since we only have two choices of where to send elements of $A$, I think its intuitive that for a given subset of $A$ we will send all elements of that subset to 1, and every other element of $A$ to 0. Writing this out mathematically we define a function $\phi: \mathscr{P}(A) \to 2^A$ by the following. Let $N \in \mathscr{P}(A)$ be arbitrary. We then have that $\phi(N) = f_N$ where

\[ f_N(a) = \begin{cases}
    1 & a \in N \\ 0 & a \in A / N
\end{cases}
\]

We now show that $\phi$ is a bijection

\textit{Injective}: Suppose that $\phi(M)=\phi(N)$. We must then have that $f_N$ and $f_M$ are equal as functions. This is only possible if they send the exact same elements to 1 and the same elements to 0, but by construction this is only possible if $M=N$, proving injectivity.

\textit{Surjective}: Let $f:A \to \{0,1\}$ be any function. We have that $f^{-1}(1) \subseteq A$ and so $f^{-1}(1) \in \mathscr{P}(A)$. We then immediately get that $\phi(f^{-1}(1)) = f$, proving surjectivity.

\qed

\bigskip

\subsection{Categories}

\medskip

\textbf{Exercise 3.1.} Let $\mathsf{C}$ be a category. Consider a structure $\mathsf{C}^{op}$ with \begin{itemize}
    \item $\text{Obj}(\mathsf{C}^{op}) := \text{Obj}(\mathsf{C})$
    \item for $A,B$ objects of $\mathsf{C}^{op}$, $\text{Hom}_{\mathsf{C}^{op}}(A,B) := \text{Hom}_{\mathsf{C}}(B,A)$
\end{itemize}
Show how to make this into a category (that is, define composition of morphisms in $\mathsf{C}^{op}$ and verify the properties listed in \S 3.1 )

\medskip

\textbf{Solution} Let A,B be objects in $\mathsf{C}^{op}$ and $f \in \text{Hom}_{\mathsf{C}^{op}}(A,B)$ be a morphism. Then $f \in \text{Hom}_{\mathsf{C}}(B,A)$. Similarly let $g \in \text{Hom}_{\mathsf{C}^{op}}(B,C)$, so $g \in \text{Hom}_{\mathsf{C}}(C,B)$. We define $gf \in \text{Hom}_{\mathsf{C}^{op}}(A,C)$ as the composition $fg \in \text{Hom}_{\mathsf{C}}(C,A)$. We no verify each property of a category holds.

The identity morphism $1_A \in \Hom_{\Cat^{op}}(A,A)$ is inherited from the original category by definition. Let $f \in \Hom_{\Cat^{op}}(A,B)$, $g \in \Hom_{\Cat^{op}}(B,C)$, $h \in \Hom_{\Cat^{op}}(C,D)$. For clarity we denote $\overline{fg}$ for a composition in the original category. We have the following:

\[ (gf)h = (\overline{fg})h = \overline{h}(\overline{fg}) = (\overline{hf})\overline{g} = g(fh) \]

And so composition is associative. The requirement for $\Hom_{\Cat^{op}}(A,B)$ and $\Hom_{\Cat^{op}}(C,D)$ being disjoint unless $A=C$ and $B=D$ is also inherited from $\Cat$ and so we have verified that $\Cat^{op}$ is a valid category.

\qed

\bigskip

\Exercise{3.2.} If $A$ is a finite set, how large is $\End_{\cat{Set}}(A)$?

\Solution We know that $\End_{\cat{Set}}(A) = \Hom(A,A) = A^A$. From Exercise (2.10) we know that $|A^A| = |A|^{|A|}$.

\qed

\bigskip

\Exercise{3.3.} Formulate precisely what it means to say that $1_a$ is an identity with respect to composition in Example 3.3, and prove this assertion.

\Solution Saying that $1_a \in \Hom(a,a)$ is the identity with respect to composition, means that for $f \in \Hom(a,b)$ we have

\[  f 1_b = (a,b), 1_a f = (a,b)  \]

We now prove this claim. If $f \in \Hom(a,b)$ then this implies that $a \sim b$. $1_a = (a,a)$ implies $a \sim a$ and similarly $1_b$ implies $b \sim b$, which we now to be true as the relation is defined to by reflexive. By the transitivity of the relation we have that

\[ a \sim b, b \sim b \implies a \sim b \]

\[ a \sim a, a \sim b \implies a \sim b \]

And this is exactly equivalent to the statement above.

\qed

\bigskip

\Exercise{3.4.} Can we define a category in the style of Example 3.3 using the relation $<$ on the set $\mathbb{Z}$?

\Solution Example 3.3 uses a relation that is both reflexive and transitive to construct a category. The relation $<$ is not reflexive, aas an integer cannot be strictly greater/lesser than itself, and so we cannot use Example 3.3 to define a category. In general, we are not going to have an identity morphism in this case, as for $a \in \mathbb{Z}$ the set $\Hom(a,a)$ will be empty for the reason above.

\qed

\Exercise{3.5.} Explain in what sense Example 3.4 is an instance of the categories considered in Example 3.3.

\Solution The power set $\mathscr{P}(S)$ of a set $S$ is itself a set, and $\subseteq$ is a relation that is both reflexive and transitive that defines the morphisms of the category. Its immediately clear that from this that Example 3.4 is a particular case of Example 3.3. 

\qed

\Exercise{3.6.} (Assuming some familiarity with linear algebra.) Define a category $\cat{V}$ by taking $\Obj(\cat{V}) = \mathbb{N}$ and letting $\Hom_V(n,m) = $ the set of $m \times n$ matrices with real entries, for all $m,n \in \mathbb{N}$. (I will leave the read the task of making sense of a matrix with 0 rows or columns.) Use product of matrices to define composition. Does this category 'feel' familiar?

\Solution We simply show that each requirement of the category $\cat{V}$ is satisfied. For every $n \in \mathbb{N}$ we can define our identity morphism as the $n-$th dimensional identity matrix ($I_n$). We clearly have that $I_n \in \Hom(n,n)$. Let $M \in \Hom(n,m), N \in \Hom(m,u), P \in \Hom(u,v)$. The the composition of these morphisms, $PNM$, is just the product of three matrices, and we know that matrix multiplication is associative so composition is associative in $\cat{V}$. We next show that $I_n$ is the identity with respect to composition. For $M \in \Hom(n,m)$ we have that

\[ M I_n = M, I_m M = M \]

And so this requirement is satisfied. Finally we know that $\Hom(n,m) = \Hom(u,v)$ if and only if $n=u$ and $m=v$ as the equality of the sets means the dimensions of the matrices within each set must be the same. We have shown that $\cat{V}$ is a well defined category.

\qed

\Exercise{3.7.} Define carefully the objects and morphisms in Example 3.7, and draw the diagram corresponding to composition.

\medskip

\Solution Defining this category is very similar to how $\Cat_A$ was defined in the example prescribed. We will denote this new category $\Cat^A$. The objects of $\Cat^A$ are all morphisms from the object $A$ of $\Cat$ to any object of $\Cat$. We define this pictorially by using a 'bottom-up' approach opposite to that of Example 3.5

\[
\begin{tikzcd}
Z \\ A \arrow[u,"f"] 
\end{tikzcd}
\]

Let $f_1,f_2$ be objects of $\Cat^A$
\[
\begin{tikzcd}
    Z_1 & Z_2 \\ A \arrow[u,"f_1"] & A \arrow[u,"f_2"]
\end{tikzcd}
\]

Then, morphisms $f_1 \to f_2$ of $\Cat^A$ are defined to be commutative diagrams

\[
\begin{tikzcd}[column sep=small]
    Z_1 \arrow[rr,"\gamma"] & & Z_2 \\ & A \arrow[ul,"f_1"] \arrow[ur,"f_2"]
\end{tikzcd}
\]

Or equivalently the morphism is exactly the morphism $\gamma: Z_1 \to Z_2$ such that $f_2 = \gamma f_1$. The composition of two morphisms $f_1 \to f_2 \to f_3$ corresponds to joining the two commutative diagrams

\[
\begin{tikzcd}
    Z_1 \arrow[r,"\gamma_1"] & Z_2 \arrow[r,"\gamma_2"] & Z_3 \\ & A \arrow[ul,"f_1"] \arrow[u,"f_2"] \arrow[ur,"f_3" below]
\end{tikzcd}
\]

The requirements for $\Cat^A$ to be a category are shown exactly the same way as in Example 3.5

\qed
\bigskip

\Exercise{3.8.} A \textit{subcategory} $\Cat'$ of $\Cat$ consists of a collection of objects in $\Cat$ with sets of morphisms $\Hom_{\Cat'}(A,B) \subseteq \Hom_{\Cat}(A,B)$ for all objects $A,B$ in $\Obj(\Cat')$, such that identities and compositions in $\Cat$ make $C'$ into a category. A subcategory is \textit{full} if  $\Hom_{\Cat'}(A,B) = \Hom_{\Cat}(A,B)$ for all $\Obj(\Cat')$. Construct a category of \textit{infinite sets} and explain how it may be viewed as a full subcategory of $\cat{Set}$.

\medskip

\Solution Let $\cat{Setinf}$ denote the category of infinite sets. Clearly, the objects of $\cat{Setinf}$ are infinite sets, and morphisms from $A$ to $B$ are the set of all functions from $A$ to $B$, $B^A$. It's (hopefully) immediately clear that all of the requirements for this to be a valid category are inherited from $\cat{Set}$, Since $A$ and $B$ are also objects of $\cat{Set}$, and from this angle its clear that $\cat{Setinf}$ is a subcategory of $\cat{Set}$. Furthermore, since $\Hom_{\cat{Setinf}}(A,B) = B^A = Hom_{\cat{Set}}(A,B)$, we can also see that this subcategory is full.

\qed

\bigskip

\Exercise{3.9.} An alternative to the notion of \textit{multiset} introduced in \S 2.2 is obtained by considering sets endowed with equivalence relations; equivalent elements are taken to be multiple instances of elements 'of the same kind'. Define a notion of morphism between such enhanced sets, obtaining a category $\cat{MSet}$ containing (a 'copy' of) $\cat{Set}$ as a full subcategory. Which objects in $\cat{Mset}$ determine ordinary multisets as defined in \S 2.2 and how? Spell out what a morphism of multisets would be from this point of view. (There are several natural notions of morphisms of multisets. Try to define morphisms in $\cat{MSet}$ so that the notion you obtain for ordinary multisets captures your intuitive understanding of these objects.) 

\Solution As mentioned by the question we will let $\Obj(\cat{MSet})$ be Sets endowed with an equivalence relation between its elements. We will define a morphism to be a function between equivalence classes between the two enhanced sets. We show that this is a well defined category. It may help first to make what we have said more explicit

\smallskip

Our objects will be 2-tuples $(A,\sim)$ where A is a set, and $\sim$ is an equivalence relation endowed on the set. For, objects $(A,\sim_A)$, $(B,\sim_B)$ we will define a morphism to be a function $f:(A,\sim_A) \to (B,\sim_B)$ such that if $x \sim_A y$ for $x,y \in A$, then $f(x) \sim_B f(y)$. Clearly there is an identity morphism $I:(A,\sim_A) \to (A,\sim_A)$ which is just the identity function mapping $x \mapsto x$. For composition of morphisms, we use usual function composition. For $f \in \Hom(A,B)$ and $g \in \Hom(B,C)$, if $x \sim_A y$ for $x,y \in A$, then $f(x) \sim_B f(y)$. We then must also have that $g(f(x)) \sim_C g(f(y))$ and so $gf \in \Hom(A,C)$. Associativity comes immediately from the fact that these are functions. It is clear that for the object $(S,=)$ that this is just a normal set, and so for another set of the same nature, say $(M,=)$ we have that the morphisms between the two objects are just regular set-functions exactly as described in the category $\cat{Set}$. Therefore, it should hopefully be clear that $\cat{Set}$ is a full subcategory of $\cat{MSet}$.

\qed

\bigskip

\Exercise{3.10.} Since the objects of a category $\Cat$ are not (necessarily interpreted as) set, it is not clear how to make sense of a notion of 'subobject' in general. In some situations it \textit{does} make sense to talk about subobjects, and the subobjects of any given object $A$ in $\Cat$ are in one-to-one correspondence with the morphisms $A \to \Omega$ for a fixed, special object $\Omega$ of $\Cat$, called a \textit{subobject classifier}. Show that $\cat{Set}$ has a subobject classifier.

\Solution The subobjects of $A$ in this case will be subsets. We know that all of the subsets of $A$ are captured by its power set $\mathscr{P}(A)$. Therefore we need only find a one to one correspondence between the power set and the morphisms to a fixed set. One may recall this is exactly what we did with $2^A$. That is, we considered functions $f:A \to \{0,1\}$ and found a one to one correspondence between these and the power set of $A$. This is exactly what we are required to show, and so Exercise (2.11) directly shows that $\{0,1\}$ is the subobject classifier of $\cat{Set}$.

\qed

\bigskip

\Exercise{3.11.} Draw the relevant diagrams and define composition and identities for the category $\Cat^{A,B}$ mentioned in Example 3.9. Do the same for the category $\Cat^{\alpha,\beta}$ mentioned in Example 3.10.

\Solution For $\Cat^{A,B}$ we have that $\Obj(\Cat^{A,B})$ are the diagrams

\[
\begin{tikzcd}[column sep=small]
    & & A \arrow[lld,"f" above] \\ Z \\ & & B \arrow[llu, "g"]
\end{tikzcd}\]

Or equivalently pairs of morphisms in $\Cat$, $(f,g)$ where $f:A \to Z$, $g:B \to Z$. The morphisms $(f_1,g_1) \to (f_2,g_2)$ are commutative diagrams

\[
\begin{tikzcd}
    & & A \arrow[lld, "f_1" above, bend right=20] \arrow[ld,"f_2"] \\ Z_1 \arrow[r, "\sigma"] & Z_2 & \\ & & B \arrow[llu,"g_1", bend left=20] \arrow[lu, "g_2"]
\end{tikzcd}\]

Or equivalently the morphism $\sigma:Z_1 \to Z_2$ such that $f_2=\sigma f_1$ AND $g_2 =\sigma g_1$. The identity morphism will just be the identity function $\text{id}:Z \to Z$ in the following diagram:

\[
\begin{tikzcd}
    & & A \arrow[lld, "f" above, bend right=20] \arrow[ld,"f"] \\ Z_1 \arrow[r, "\text{id}"] & Z_2 & \\ & & B \arrow[llu,"g", bend left=20] \arrow[lu, "g"]
\end{tikzcd}\]

as $f = f \id$ and $g = g \id$. For composition, this is just joining two diagrams together and this process is independent of the order join the diagrams together, so it is associative.

The process for $\Cat^{\alpha,\beta}$ is very similar. $\Obj(\Cat^{\alpha,\beta})$ are the commutative diagrams

\[
\begin{tikzcd}
    & A \arrow[ld, "f"] & \\ Z & & C \arrow[lu, "\alpha"] \arrow[ld,"\beta"] \\ & B \arrow[lu, "g"] &
\end{tikzcd}\]


or equivalently paris of morphisms $(f,g)$ Such that $f\alpha=g\beta$. Then, morphisms correspond to the commutative diagram
\[
\begin{tikzcd}
    & & A \arrow[lld, "f_1" above, bend right=20] \arrow[ld,"f_2"] & \\ Z_1 \arrow[r, "\sigma"] & Z_2 & & C \arrow[lu,"\alpha"] \arrow[ld,"\beta"] \\ & & B \arrow[llu,"g_1", bend left=20] \arrow[lu, "g_2"] &
\end{tikzcd}\]

With the identity being the same as the previous category, and associativity having an analogous argument.

\qed

\end{document}